\documentclass[10pt]{article}
\usepackage{amsmath, amssymb, amsthm}
\usepackage{booktabs}
\usepackage{rotating}

\begin{document}
\noindent\textbf{\Large Math walkthrough for $F(m,n)$ when $n = 2^k$}

\section{FLOP Count for each step}

See the table.

\begin{sidewaystable}
    \centering
    \label{tab:my_label}
  \begin{tabular}{llcc}
    \toprule
    Operation & Function & FLOPS($n_1,n_2,m_3$) & FLOPS($n_1=n_2=n, m_3=m-2n$)\\
    \midrule
    Factor left panel & GEQRT3($m\times n_1$, $n_1\times n_1$) & $F(m,n_1)$ & $F(m,n)$\\
    \midrule
    Update trailing matrix & TRMM(L, $n_1\times n_1$, $n_1\times n_2$) & $n_1^2n_2$ & $n^3$\\
    & GEMM($n_1\times n_2$, $n_2\times n_2$, $n_1\times n_2$) & $2n_1n_2n_2$ & $2n^3$\\
    & GEMM($n_1\times m_3$, $m_3 \times n_2$, $n_1\times n_2$) & $2n_1n_2m_3$ & $2n^2(m-2n)$\\
    & TRMM(L, $n_1\times n_1$, $n_1\times n_2$) & $n_1^2n_2$ & $n^3$\\
    & GEMM($n_2\times n_1$, $n_1\times n_2$, $n_2\times n_2$) & $2n_2n_2n_1$ & $2n^3$\\
    & GEMM($m_3\times n_1$, $n_1\times n_2$, $m_3\times n_2$) & $2m_3n_2n_1$ & $2(m-2n)n^2$\\
    & TRMM(L,$n_1\times n_1$, $n_1\times n_2$) & $n_1^2n_2$ & $n^3$\\
    & DIFF($n_1\times n_2$, $n_1 \times n_2$) & $n_1n_2$ & $n^2$\\
    \midrule
    Factor right panel & GEQRT3($(m - n_1)\times n_2$, $n_2\times n_2$) &
    $F(m-n_1,n_2)$ & $F(m-n,n)$ \\
    \midrule
    Construct T & TRMM(R, $n_2\times n_2$, $n_1\times n_2$) & $n_1n_2^2$ & $n^3$\\
    & GEMM($n_1\times m_3$, $m_3\times n_2$, $n_1\times n_2$) & $2n_1n_2m_3$ & $2n^2(m-2n)$\\
    & TRMM(L, $n_1\times n_1$, $n_1\times n_2$) & $n_1^2n_2$ & $n^3$\\
    & TRMM(R, $n_2\times n_2$, $n_1\times n_2$) & $n_1n_2^2$ & $n^3$\\
    \bottomrule
  \end{tabular}
\end{sidewaystable}



\section{Recurrence from the code}

For $n == 1$ the code does one $larfg$ on a vector of length $m$, giving
\[F(m,1) = 3 m\]
For $n > 1$ and $n$ a power of two, the code splits
\[n_1 = n/2\]
\[n_2 = n/2\]
and the recurrence becomes
\[F(m,n) = F(m,n/2) + F(m-n/2,n/2) + C(m,n)\]
where the local flop contribution from the helper calls is
\[C(m,n) = n_1 n_2 (6 m - 2 n_1 + 1)\]
For $n_1=n_2=n/2$ this is
\[C(m,n) = \frac{n^2}{4}  (6 m - n + 1)\]

\section{Key simplification: $F(m,n)$ is linear in $m$}

We claim that for $n=2^k$, the recurrence
\[ F(m,n) = F\!\left(m,\frac{n}{2}\right) + F\!\left(m-\frac{n}{2},\frac{n}{2}\right) + \frac{n^2}{4}(6m-n+1) \]
with base case
\[ F(m,1)=3m \]
has the form
\[ F(m,n)=3n^2m+b(n), \]
where $b(n)$ is independent of $m$.
\begin{proof}
We proceed by induction on $n$.

For the base case $n=1$,
\[ F(m,1)=3m, \]
which can be written as
\[ F(m,1)=3(1)^2m+b(1) \]
with $b(1)=0$. Thus the claim holds for $n=1$.

Now assume that for some power of two $n/2$,
\[ F(m,n/2) = 3\left(\frac{n}{2}\right)^2m + b\!\left(\frac{n}{2}\right).  \]
This is the induction hypothesis.

Substituting the induction hypothesis into the recurrence gives
\begin{align*}
F(m,n) &= \left[ 3\left(\frac{n}{2}\right)^2m + b\!\left(\frac{n}{2}\right) \right] \\
&\quad+ \left[ 3\left(\frac{n}{2}\right)^2 \left(m-\frac{n}{2}\right) + b\!\left(\frac{n}{2}\right) \right] \\
&\quad+ \frac{n^2}{4}(6m-n+1).
\end{align*}

Collecting the terms containing $m$,
\begin{align*}
\operatorname{coeff}_m(F(m,n))
&= 3\left(\frac{n}{2}\right)^2 + 3\left(\frac{n}{2}\right)^2 + \frac{n^2}{4}\cdot 6 \\
&= \frac{3n^2}{4} + \frac{3n^2}{4} + \frac{3n^2}{2} \\
&= 3n^2.
\end{align*}

Hence
\[
F(m,n) = 3n^2m + \Biggl( 2b\!\left(\frac{n}{2}\right) -\frac{3n^3}{8} -\frac{n^3}{4} +\frac{n^2}{4} \Biggr).
\]

The expression in parentheses depends only on $n$, so we define
\[ b(n) = 2b\!\left(\frac{n}{2}\right) -\frac{5}{8}n^3 +\frac{1}{4}n^2.  \]

Therefore
\[ F(m,n)=3n^2m+b(n), \]
which completes the induction.
\end{proof}

%Look at the $m$-dependence:
%\begin{enumerate}
    %\item $F(m,n/2)$ has slope $3 (n/2)^2 = 3 n^2 / 4$
    %\item $F(m-n/2,n/2)$ also has slope $3 n^2 / 4$
    %\item $C(m,n)$ has slope $6 n_1 n_2 = 6 (n/2)^2 = 3 n^2 / 2$
%\end{enumerate}
%So the total slope in $m$ is
%\[3 n^2 / 4 + 3 n^2 / 4 + 3 n^2 / 2 = 3 n^2\]
That means for each fixed power-of-two $n$ we can write
\[F(m,n) = 3 n^2 m + b(n)\]
with $b(n)$ independent of $m$.


\section{Recurrence for the offset $b(n)$}

Plugging $F(m,n)=3 n^2 m + b(n)$ back into the recurrence gives
\[b(n) = 2 b(n/2) - (5/8) n^3 + (1/4) n^2\]
This is valid for $n = 2^k$ with $b(1)=0$.


\section{Solve $b(n)$ for $n = 2^k$}

Define $n = 2^k$ and $g(k) = b(2^k)$.
Then the recurrence becomes
\[g(k) = 2 g(k-1) - (5/8) * 2^{3k} + (1/4) * 2^{2k}\]
Divide by $2^k$ and set $h(k) = g(k)/2^k$:
\[h(k) = h(k-1) - (5/8) 4^k + (1/4) 2^k\]
with $h(0)=0$.
Now sum over $j$ where $1\leq j \leq k$:
\[h(k) = - (5/8) \sum_{j=1}^k 4^j + (1/4) \sum_{j=1}^k 2^j\]
Use the geometric sums:
\[ \sum_{j=1}^k 4^j = (4^{k+1}-4)/3\]
\[ \sum_{j=1}^k 2^j = 2^{k+1}-2\]
So
\[h(k) = \left(-\frac{5}{8}\right) \left(\frac{4^{k+1}-4}{3}\right) +
\left(\frac{1}{4}\right)\left(2^{k+1}-2\right)\]
Because $n=2^k$, this is
\[h(k) = -\frac{5}{24}(4 n^2 - 4) + \frac{1}{4}(2n - 2)\]
Simplify carefully:
\[h(k) = -\frac{5}{6} n^2 + \frac{1}{2} n + \frac{1}{3}\]
Then
\[b(n) = n h(k) = -\frac{5}{6} n^3 + \frac{1}{2} n^2 + \frac{1}{3} n\]


\section{Final closed form}

For $n = 2^k$, the flop count is
\[F(m,n) = 3 m n^2 - \frac{5}{6} n^3 + \frac{1}{2} n^2 + \frac{1}{3} n\]


\section{Base-case sanity check}

For $n = 1$ this gives:
\[F(m,1) = 3m - 5/6 + 1/2 + 1/3 = 3m\]
so the formula is consistent with the base case.


\section{Numeric verification}

I also verified this formula against the exact recursive count for powers of two:
\begin{enumerate}
    \item $n = 2, 4, 8, 16, 32, 64$
    \item and several $m$ values ($m=n$, $m=n+1$, $m=2n$, etc.)
\end{enumerate}
All matched exactly.


\section{Summary}

So the clean closed form for powers-of-two $n$ is:
\[F(m,n) = 3 m n^2 - \frac{5}{6} n^3 + \frac{1}{2} n^2 + \frac{1}{3} n\]
\end{document}
